\section{Petunjuk Penggunaan Buku Ajar}

\subsection{Untuk Mahasiswa}

Buku ini dirancang untuk dipelajari secara bertahap selama 16 pertemuan (2 SKS). Sebelum perkuliahan, bacalah Sub-CPMK di awal setiap bab untuk memahami target pembelajaran, pelajari materi pokok dengan seksama, dan jalankan semua contoh kode menggunakan browser dan editor teks seperti VS Code \cite{vscode}. Mencatat pertanyaan atau konsep yang belum dipahami akan memperkaya diskusi di kelas.

Persiapkan lingkungan development sebelum memulai. Instal XAMPP untuk menjalankan PHP dan MySQL secara lokal, serta VS Code dengan addon untuk HTML, CSS, JavaScript, dan PHP \cite{xampp}. Buka folder proyek di VS Code dan praktikkan setiap contoh kode dengan membuat file HTML, CSS, JS, atau PHP baru. Jalankan file PHP melalui \texttt{http://localhost/nama-folder/file.php} setelah Apache dan MySQL aktif di XAMPP Control Panel.

Selama perkuliahan, manfaatkan waktu untuk diskusi, praktikum, dan live coding. Kerjakan aktivitas pembelajaran secara aktif, tanyakan hal-hal yang belum jelas, dan berpartisipasi dalam peer review terhadap tampilan atau kode web yang dibuat rekan. Setelah perkuliahan, kerjakan latihan dan refleksi, lakukan asesmen mandiri, centang checklist kompetensi yang telah dikuasai, dan kerjakan proyek mini untuk memperdalam pemahaman \cite{webdev}.

\subsection{Untuk Dosen}

Buku ini dapat digunakan sebagai bahan ajar utama, sumber latihan dan tugas, referensi untuk menyusun soal ujian, dan panduan untuk merancang aktivitas pembelajaran. Dokumentasi daring seperti MDN Web Docs, W3Schools, PHP Manual, dan MySQL Documentation dapat dilengkapi sebagai referensi tambahan sesuai kebutuhan \cite{mdn-learn}.

Rekomendasi penggunaan: gunakan bagian Sub-CPMK untuk menyusun indikator penilaian; bagian Aktivitas untuk tugas praktikum; bagian Latihan untuk kuis atau PR; bagian Asesmen untuk rubrik proyek; dan bagian Checklist untuk self-assessment mahasiswa. Komposisi 70\% praktek dan 30\% teori tercermin dalam porsi contoh kode dan latihan coding yang banyak di setiap bab.
